\documentclass[]{article}

\begin{document}
%The abstract must be between 200 and 1000 characters.
This PhD project develops a human–AI framework for bridge structural health monitoring. The framework leverages multi-modal data, structural information, inspection histories, environmental observations, and engineering knowledge through a digital twin. Evidence is semantically contextualized into an evolving representation of the bridge condition. AI supports the interpretation of this evidence by identifying plausible deterioration mechanisms and assessing maintenance plans. Extended Reality creates an immersive environment to support users' decision making. The framework is validated experimentally, considering two bridge case studies in the Global South (Brazil and Kenya). Lorenzo Loyola (PhD candidate) is currently employed at Aarhus University (AU). The project establishes a collaboration between AU and the research group of Prof. Ueli Angst at ETH Zurich.
\end{document}